% Part: first-order-logic
% Chapter: model-theory
% Section: theory-of-m

\documentclass[../../../include/open-logic-section]{subfiles}

\begin{document}

\olfileid{mod}{bas}{thm}

\section{\printtoken{S}{structure} ची उपपत्ती}

प्रत्येक !!{structure}~$\Struct{M}$ काही !!{sentence}s सत्य ठरवते आणि
काही असत्य ठरवते. ती सत्य ठरवत असलेल्या सर्व !!{sentence}s च्या संचाला
तिची \emph{उपपत्ती} म्हणतात. तो संच खरोखरच एक उपपत्ती असतो, कारण
त्याच्यापासून तार्किकरीत्या निष्पन्न होणारे काहीही त्याच्या सर्व
प्रतिमानांमध्ये, आणि त्यांत~$\Struct{M}$ चाही समावेश होतो, सत्य असले
पाहिजे.

\begin{defn}
  !!a{structure}~$\Struct M$ दिलेली असताना, $\Struct{M}$ ची
  \emph{उपपत्ती} म्हणजे $\Struct{M}$ मध्ये सत्य असलेल्या !!{sentence}s
  चा संच $\Theory{M}$ होय; म्हणजे, $\Theory{M} =
  \Setabs{!A}{\Sat{M}{!A}}$.
\end{defn}

अभिप्रेत अर्थनिर्धारण असलेल्या !!{sentence}s च्या संचांसाठीही आपण
``उपपत्ती'' हा शब्द अनौपचारिकपणे वापरतो, मग ते संच निगामीदृष्ट्या बंद
असोत किंवा नसोत.

\begin{prop}
कोणत्याही $\Struct{M}$ साठी $\Theory{M}$ संपूर्ण असते.
\end{prop}

\begin{proof}
कोणत्याही !!{sentence}~$!A$ साठी एकतर $\Sat{M}{!A}$ किंवा
$\Sat{M}{\lnot !A}$; म्हणून एकतर $!A \in \Theory{M}$ किंवा
$\lnot !A \in \Theory{M}$.
\end{proof}

\begin{prop}\ollabel{prop:equiv}
  प्रत्येक $!A \in \Theory{M}$ साठी $\Struct{N} \models !A$ असेल,
  तर $\Struct{M} \elemequiv \Struct{N}$.
\end{prop}

\begin{proof}
सर्व $!A \in \Theory{M}$ साठी $\Sat{N}{!A}$ असल्यामुळे
$\Theory{M} \subseteq \Theory{N}$. जर $\Sat{N}{!A}$, तर
$\Sat/{N}{\lnot !A}$; म्हणून $\lnot !A \notin \Theory{M}$.
$\Theory{M}$ संपूर्ण असल्यामुळे $!A \in \Theory{M}$. म्हणून
$\Theory{N} \subseteq \Theory{M}$ आणि आपल्याला
$\Struct{M} \elemequiv \Struct{N}$ मिळते.
\end{proof}

\begin{rem}\ollabel{remark:R}
  $\Struct{R} = \langle\Real, <\rangle$ विचारात घ्या. या
  !!{structure} चा प्रांत म्हणजे वास्तव संख्यांचा संच~$\Real$ असून,
  ती फक्त एक 2-स्थानी !!{predicate} असलेल्या !!{language} साठीची रचना
  आहे; त्या विधेयाचे अर्थनिर्धारण वास्तव संख्यांवरील $<$ संबंध असे केले
  आहे. स्पष्टपणे $\Struct{R}$ ही !!{nonenumerable} आहे; पण
  $\Theory{R}$ उघडपणे सुसंगत असल्यामुळे लोव्हेनहाइम--स्कोलेम प्रमेयाने
  तिला एक !!a{enumerable} प्रतिमान असते, त्याला $\Struct{S}$ म्हणू.
  \olref{prop:equiv} नुसार $\Struct{R} \equiv \Struct{S}$. शिवाय,
  $\Struct{R}$ आणि $\Struct{S}$ या समरूपी नसल्यामुळे यावरून
  \olref[iso]{thm:isom} चा उलट दावा सर्वसाधारणपणे खोटा ठरतो.
\end{rem}

\end{document}
